\documentclass[11pt]{article}

\usepackage{amsmath,amssymb}
\usepackage{xurl}
\usepackage[colorlinks=true,linkcolor=blue,citecolor=blue,urlcolor=blue]{hyperref}

\input{./command.tex}

\title{Yamagami Lemma~6: The Strict-Disk Endpoint at Stride One}
\author{}
\date{2026-08-24}

\begin{document}
\maketitle
\gapnote{false-source}{resolved}

\section{The source statement}

Let $C$ be the permutation matrix for the cyclic shift on
$\mathbb Z/n\mathbb Z$. For integers $l,m,n,s$ in the range
\begin{align}
    1
        &\leq m
        \leq n-l,
    \label{eq:yamagami93_lemma6_endpoint_parameter_range_m}\\
    n
        &\leq s,
    \label{eq:yamagami93_lemma6_endpoint_parameter_range_s}
\end{align}
Lemma~6 of Ref.~\cite{Yamagami1993Cyclic} considers the cyclic matrix
\begin{align}
    S_{n,m,s}
        &=
        (s-m)I+C^l+C^{l+1}+\cdots+C^{l+m-1}.
    \label{eq:yamagami93_lemma6_endpoint_cyclic_matrix}
\end{align}
Its row sum is $s$. The lemma states that every eigenvalue $\lambda\neq s$
lies in the open disc
\begin{align}
    \left|\lambda-\frac{s}{2}\right|
        &< \frac{s}{2}.
    \label{eq:yamagami93_lemma6_endpoint_open_disk}
\end{align}
This is the spectral condition used by Corollary~5 of
Ref.~\cite[pp.~523--524]{Yamagami1993Cyclic}.

\section{The false stride-one endpoint}

Choose the stride-one endpoint
\begin{align}
    l
        &= 1,
    \label{eq:yamagami93_lemma6_endpoint_stride_one}\\
    m
        &= n-1,
    \label{eq:yamagami93_lemma6_endpoint_maximal_length}\\
    s
        &= n.
    \label{eq:yamagami93_lemma6_endpoint_minimal_row_sum}
\end{align}
These values satisfy
(\ref{eq:yamagami93_lemma6_endpoint_parameter_range_m}) and
(\ref{eq:yamagami93_lemma6_endpoint_parameter_range_s}). The matrix in
(\ref{eq:yamagami93_lemma6_endpoint_cyclic_matrix}) becomes
\begin{align}
    S_{n,n-1,n}
        &=
        I+C+C^2+\cdots+C^{n-1}.
    \label{eq:yamagami93_lemma6_endpoint_character_sum_matrix}
\end{align}
Every nontrivial standard character of $\mathbb Z/n\mathbb Z$ is an
eigenvector of
(\ref{eq:yamagami93_lemma6_endpoint_character_sum_matrix}) with eigenvalue
zero. For such an eigenvalue,
\begin{align}
    \left|0-\frac{n}{2}\right|
        &= \frac{n}{2}.
    \label{eq:yamagami93_lemma6_endpoint_boundary_equality}
\end{align}
Equation~(\ref{eq:yamagami93_lemma6_endpoint_boundary_equality}) contradicts
the strict inequality in
(\ref{eq:yamagami93_lemma6_endpoint_open_disk}). At this endpoint, the strict
triangle-inequality step in the proof of Lemma~6 of
Ref.~\cite{Yamagami1993Cyclic} also loses all strictness.

\section{The correction used by QICLean}

The Choi-type application requires only stride one with $m\leq n-2$. Under
the assumptions
\begin{align}
    n
        &\geq 3,
    \label{eq:yamagami93_lemma6_endpoint_corrected_n}\\
    s
        &\geq n,
    \label{eq:yamagami93_lemma6_endpoint_corrected_s}\\
    m
        &\leq n-2,
    \label{eq:yamagami93_lemma6_endpoint_corrected_m}
\end{align}
the formal development proves
(\ref{eq:yamagami93_lemma6_endpoint_open_disk}) for every nonzero standard
character.\footnote{Formal declaration:
\leanid{Yamagami.norm_symbol_sub_half_lt_half} in
\doclink{QICLean/Analysis/CyclicReciprocal}.}
The proof follows the two cases used at $s=n$ in Lemma~6 of
Ref.~\cite{Yamagami1993Cyclic} and then applies the source's translation
argument from $n$ to $s$.

The omitted endpoint $m=n-1$ also satisfies the open-disc condition when
$s>n$. Its eigenvalue on every nontrivial character is
\begin{align}
    \lambda
        &= s-n.
    \label{eq:yamagami93_lemma6_endpoint_shifted_endpoint_eigenvalue}
\end{align}
For this eigenvalue, the open-disc condition is equivalent to
\begin{align}
    \left|s-n-\frac{s}{2}\right|
        &< \frac{s}{2},
    \label{eq:yamagami93_lemma6_endpoint_shifted_endpoint_disk}\\
    0
        &< n<s.
    \label{eq:yamagami93_lemma6_endpoint_shifted_endpoint_condition}
\end{align}
Thus the corrected stride-one alternatives are
\begin{align}
    m
        &\leq n-2
        &&\text{with }s\geq n,
    \label{eq:yamagami93_lemma6_endpoint_corrected_main_range}\\
    m
        &= n-1
        &&\text{with }s>n.
    \label{eq:yamagami93_lemma6_endpoint_corrected_endpoint_range}
\end{align}

\section{Scope}

\textbf{Verdict: false source endpoint, resolved for the required stride-one
range.} This correction concerns only the open-disc condition in Lemma~6 of
Ref.~\cite{Yamagami1993Cyclic}; it does not by itself prove Yamagami's cyclic
reciprocal inequality.

The finite Nowosad theorem~\cite{Nowosad1968Isoperimetric}, the uniqueness
assertion in Lemma~3 of Ref.~\cite{Yamagami1993Cyclic}, and the
simultaneous-singularity boundary count are now formalized. The cyclic matrix
package in \doclink{QICLean/Analysis/YamagamiCyclicMatrix} proves
invertibility, positive semidefiniteness of the Hessian with kernel
$\mathbb R\mathbf 1$, and the bridge from the generic Nowosad functional to the
concrete cyclic functional \leanid{Yamagami.functional}. The regular-boundary
recurrence on p.~525 of Ref.~\cite{Yamagami1993Cyclic} is formalized by
\leanid{Yamagami.functional_le_deleteLastCoordinate}.

The declaration \leanid{Yamagami.functional_card_le_one} assembles these
ingredients. Its compact-superlevel argument produces a positive maximizer,
and the Fourier--Hessian and Nowosad uniqueness theorem makes that maximizer
scalar. The endpoint correction is therefore a resolved prerequisite of the
completed positivity theorem.

\bibliographystyle{alpha}
\bibliography{references}

\end{document}
